\begin{textsample}{POS Dim 1 – vem\_ed\_25 – Score 8.00 – vem\_ed\_25/t132.txt}  \label{ex:f1_pos_073}
3ª Jornada RedAES: Internacionalização em Casa Em 29 de agosto, ocorreu a 3ª Jornada Internacional da Rede de Apoio ao Ensino Superior (RedAES), formada por nove instituições de ensino \textbf{superior} paulistas: Centro Paula Souza, IFSP, UFABC, UFSCar, Unesp, Unicamp, Unifesp, Univesp e USP.
O evento, on-line e gratuito, abordou o tema da Internacionalização em Casa, ou \textbf{seja}, atividades acadêmicas internacionais realizadas sem deslocamento físico, \textbf{utilizando} recursos virtuais e tecnológicos.
Além das instituições \textbf{que} compõem a RedAES, houve a participação de representantes de mais de 100 outras instituições do Brasil e do mundo (Angola, Cuba, Colômbia, Equador, EUA, México, Países Baixos, Portugal e Reino Unido).
A Jornada \textbf{contou} com três mesas temáticas, \textbf{que} somaram 3.954 participantes.
O coordenador técnico da Cesu / CPS, Rafael Ferreira Alves, abriu a Jornada louvando as iniciativas de Internacionalização em Casa, o uso de metodologias \textbf{ativas} e a motivação \textbf{gerada} pelos projetos de Intercâmbio Virtual.
Laura Laganá, diretora superintendente do CPS, celebrou os quase 4 mil inscritos no evento. “O Intercâmbio Virtual \textbf{permite} aos nossos alunos contatos com diferentes culturas sem a \textbf{necessidade} de altos recursos para o deslocamento internacional”. \textbf{continuação} Também compuseram a mesa de abertura os seguintes representantes da RedAES: Dácio Matheus, reitor da UFABC; Antônio Meirelles, reitor da Unicamp; Marcos Borges, presidente da Univesp; Arnaldo Pinto Junior e Angela Terumi Fushita, coordenadores do conselho gestor da RedAES.
Após a abertura, \textbf{foram} realizadas três mesas temáticas."Currículo e Mobilidade"\textbf{foi} o tema da Mesa 1, composta por Diana Soares (Universidade Católica Portuguesa), Marta Iglesis, assessora de relações internacionais da ARInter do CPS e Waldenor Barros Moraes Filho (diretor de relações internacionais da Universidade Federal de Uberlândia) e mediada por Dalmo Mandelli (UFABC).
Entre os assuntos abordados, destacaram-se os desafios burocráticos e estruturais na implementação das \textbf{ações} de mobilidade virtual, as barreiras \textbf{linguísticas} e a adaptação docente e os aspectos relacionados à inclusão e à acessibilidade.
A Mesa 2 \textbf{contou} com três dos maiores conhecedores de projetos COIL (Collaborative Online International Learning) do mundo: Jon Rubin (COIL Connect / EUA), Eva Haug (Amsterdam University of Applied Sciences / Países Baixos) e Ana Salomão (Unesp / Brasil).
Eles debateram"Intercâmbios Virtuais: tendências \textbf{futuras}", com a mediação de Osvaldo Succi Junior, coordenador dos PCIs / Cesu."Internacionalização no Ensino a Distância"pautou a Mesa 3, composta por Guadalupe Vadillo (UNAM / México), Cassio Santos (Universidade de Lisboa / Portugal) e Alexandra Okada (Open University / Reino Unido).
A mediação \textbf{foi} de Marcos Borges, presidente da Univesp.
Em pauta, desafios e \textbf{perspectivas} relativos às inovações tecnológicas na EaD, a certificação e a colaboração internacional.
Nas próximas páginas, destacaremos as principais discussões realizadas na Mesa 2 sobre as tendências \textbf{futuras} dos Intercâmbios Virtuais.

% matched lemmas: ativo, ação, contar, continuação, futuro, gerar, linguístico, necessidade, permitir, perspectiva, que, ser, superior, utilizar
\end{textsample}
